% =====================================================================
\section{Шаговый двигатель}
% =====================================================================
\lessonintro
  {понять принцип шагового двигателя, подключить его через драйвер и поворачивать вал на заданный угол}
  {урок 16 (нагрузка через драйвер, отдельное питание), урок 21 (шаги по времени без \code{delay()}), урок 22 (биты)}
  {вал двигателя поворачивается на угол, введённый в Serial Monitor, и останавливается точно в нужном месте}
  {шаговый двигатель 28BYJ-48 с платой драйвера ULN2003, источник \SI{5}{V}}

\subsection{Сначала: мотор, который умеет считать}

Обычный мотор из урока 16 просто крутится, и программа не знает, на сколько он повернулся. Сервопривод
из урока 4 знает свой угол, но поворачивается только на пол-оборота. Шаговый двигатель сочетает оба
свойства: он вращается сколько угодно, но делает это маленькими одинаковыми шагами. Сосчитав шаги,
программа точно знает положение вала. Так работают 3D-принтеры, станки и механизмы камер.

\begin{center}\small
\begin{tabularx}{\textwidth}{@{}l X X X@{}}
\toprule
 & Коллекторный мотор & Сервопривод & Шаговый двигатель \\
\midrule
Движение & непрерывное вращение & поворот на угол 0--180° & вращение шагами \\
Что задаёт программа & скорость и направление & угол & число шагов и их темп \\
Знает ли положение & нет & да (внутренний датчик) & да, если не пропускал шаги \\
Где применяют & колёса, вентиляторы & рули, захваты & принтеры, станки, поворотные столы \\
\bottomrule
\end{tabularx}
\end{center}

\subsection{Как он устроен}

Внутри двигателя --- несколько обмоток (электромагнитов) вокруг ротора с магнитом. Если включать обмотки
по очереди, ротор каждый раз поворачивается к ближайшей включённой обмотке. Один такой поворот ---
шаг. Порядок включения определяет направление, темп --- скорость.

У двигателя 28BYJ-48 четыре обмотки. Если включать то одну обмотку, то две соседние, шаг получается
вдвое мельче. Этот режим называется полушаговым, и в нём восемь фаз.

\begin{figure}[H]
\centering
\begin{tikzpicture}[x=1.15cm,y=0.62cm,font=\sffamily\scriptsize]
  \foreach \row/\name/\bits in {0/IN1/{1,1,0,0,0,0,0,1},1/IN2/{0,1,1,1,0,0,0,0},2/IN3/{0,0,0,1,1,1,0,0},3/IN4/{0,0,0,0,0,1,1,1}}{
    \node[left] at (0,-\row*1.5+0.4) {\name};
    \foreach \b [count=\i from 0] in \bits{
      \ifnum\b=1
        \fill[accent!55] (\i,-\row*1.5) rectangle ++(1,0.8);
      \fi
      \draw[gray!50] (\i,-\row*1.5) rectangle ++(1,0.8);
    }
  }
  \foreach \i in {0,...,7} \node at (\i+0.5,1.3) {фаза \i};
  \draw[->] (0,-5.2) -- (8.4,-5.2) node[right]{время};
  \node[anchor=west,align=left] at (8.6,-1.9) {закрашено ---\\ обмотка включена};
\end{tikzpicture}
\caption{Полушаговый режим: восемь фаз. Включена то одна обмотка, то две соседние; после фазы 7 снова идёт фаза 0}
\end{figure}

Между валом и ротором 28BYJ-48 стоит редуктор примерно 1:64, поэтому на один оборот вала приходится
около 4096 полушагов --- меньше десятой доли градуса на шаг. Двигатель получается медленным (до
10--15 оборотов в минуту), но точным.

\subsection{Подключение}

Обмотка потребляет сотни миллиампер, поэтому между платой и двигателем стоит драйвер ULN2003 --- семь
транзисторных ключей с защитными диодами в одном корпусе (тот же ключ, что в уроке 16).

\begin{figure}[H]
\centering
\begin{tikzpicture}[font=\sffamily\scriptsize]
  \node[blk,fill=blkcpu,minimum width=2.6cm,minimum height=2.8cm] (m) at (0,0) {ESP32-S3};
  \node[blk,fill=blkper,minimum width=2.6cm,minimum height=2.8cm] (d) at (5.6,0) {Плата\\ ULN2003};
  \node[blk,fill=blkper,minimum width=2.6cm,minimum height=1.6cm] (s) at (11.6,0.6) {28BYJ-48};
  \node[blk,fill=blkrtc,minimum width=2.6cm] (p) at (11.6,-1.3) {Источник 5 В};
  \foreach \i/\pin in {1/15,2/16,3/17,4/18}
    \draw[arr,accent] ($(m.east)+(0,1.5-\i*0.6)$) -- node[above,pos=0.5]{GPIO\pin\ → IN\i} ($(d.west)+(0,1.5-\i*0.6)$);
  \draw[darr] ($(d.east)+(0,0.6)$) -- node[above]{разъём 5 проводов} (s.west);
  \draw[thick,esp] (p.west) -- node[above,text=black]{питание 5 В} ($(d.east)+(0,-1.3)$);
  \draw[thick,gray] (m.south) -- ++(0,-0.6) -| node[pos=0.25,below]{общая земля} (p.south);
\end{tikzpicture}
\caption{Подключение шагового двигателя. Сигналы идут от платы, ток обмоток --- от отдельного источника}
\end{figure}

\begin{caution}
Не питайте двигатель от вывода 3V3. От вывода 5V платы один 28BYJ-48 обычно работает, но при слабом
USB-порте плата может перезагружаться. Надёжнее отдельный источник \SI{5}{V} с общей землёй.
\end{caution}

\subsection{Программа}

Восемь фаз удобно хранить как восемь чисел: бит 0 отвечает за IN1, бит 1 --- за IN2 и так далее. Это
в точности таблица с рисунка, записанная по столбцам.

Строим программу в три шага. Подключите ULN2003 к GPIO15--GPIO18, как на рисунке.

\subsubsection{Шаг 1. Вал крутится}

Самое простое: по очереди включаем фазы 0, 1, \ldots, 7 и снова 0. Вал вращается без остановки.
Пауза между шагами отсчитывается по \code{micros()}, без \code{delay()}.

\codecap{Шаг 1. Вращение по таблице фаз}
\codeinput{l25_stepper_s1}

\begin{itemize}[leftmargin=*]
  \item \code{PHASES[]} --- таблица с рисунка: бит~0 включает IN1, бит~1 --- IN2 и так далее;
  \item \code{applyPhase()} разбирает число по битам (урок 22) и выставляет четыре вывода;
  \item \code{(phase + 1) \% 8} --- после фазы 7 снова идёт фаза 0.
\end{itemize}

\subsubsection{Шаг 2. Считаем шаги и останавливаемся}

Вместо бесконечного вращения заводим два числа: \code{current} --- где вал сейчас, \code{target} ---
куда ему нужно. Функция \code{stepperRun()} устроена как автомат из урока 21: за один вызов она делает
не больше одного шага и сразу возвращается. Пока числа различаются, двигатель шагает.
Здесь цель --- один оборот, \code{STEPS\_PER\_REV} шагов.

\codecap{Шаг 2. Один оборот и остановка}
\codeinput{l25_stepper_s2}

\begin{itemize}[leftmargin=*]
  \item \code{current += (target > current) ? 1 : -1} --- шаг в нужную сторону, поэтому вал умеет и назад;
  \item \code{((current \% 8) + 8) \% 8} --- номер фазы, который не уходит в минус при вращении назад;
  \item \code{release()} снимает ток с обмоток, когда цель достигнута, чтобы двигатель не грелся.
\end{itemize}

\subsubsection{Шаг 3. Угол из Serial Monitor}

Остаётся задавать цель во время работы. Откройте Serial Monitor и введите угол, например \code{90} или
\code{-180}: цель сдвинется на соответствующее число шагов. Вал повернётся, а программа за это время
остаётся свободной.

\codecap{Шаг 3. Поворот вала на заданный угол без \code{delay()}}
\codeinput{l25_stepper}

\begin{itemize}[leftmargin=*]
  \item \code{Serial.parseInt()} читает число из введённой строки; остаток строки отбрасываем;
  \item \code{degrees * STEPS\_PER\_REV / 360} --- перевод градусов в шаги;
  \item \code{target = current + \ldots} --- поворачиваем \emph{относительно} текущего положения.
\end{itemize}

\begin{tip}
Если шаги идут слишком часто, ротор не успевает за магнитным полем: двигатель жужжит, но не вращается,
а программа продолжает считать шаги и теряет положение. Уменьшайте \code{STEP\_US} постепенно и найдите
значение, при котором вал ещё уверенно вращается под нагрузкой.
\end{tip}

\subsection{Самостоятельный опыт: стрелочный указатель}

Наденьте на вал бумажную стрелку и нарисуйте шкалу. Потенциометр (урок 6) задаёт положение стрелки от
0° до 270°. Двигатель не знает, где он находится при включении, поэтому при запуске поставьте стрелку
на ноль вручную или добавьте кнопку «установить ноль».

\subsection{То же на ESPUI: поворот на заданный угол с телефона}

Вместо ввода угла в Serial Monitor сделаем пульт: поле «Шаг поворота» в градусах и три кнопки,
«вправо», «влево» и «в начало». Нажатия накапливаются, а шаги выполняются без \code{delay()}, как
в основной программе урока.

Телефон подключается к сети \code{S3-Stepper} (пароль \code{lesson-39-demo}), адрес панели ---
\code{http://192.168.4.1/}. Принцип тот же, что в уроке 5: обработчики панели только записывают
желание в \code{std::atomic}, а мотором управляет \code{loop()}.

\subsubsection{Шаг 1. Панель считает градусы}

Пока без мотора: поле «Шаг поворота» и кнопка «Вправо». Нажатия накапливаются в \code{requestDeg},
а \code{loop()} забирает их и прибавляет к цели, которую показывает подпись.

\codecap{Шаг 1. Number, кнопка и подпись с целью}
\codeinput[\footnotesize]{l25_espui_stepper_s1}

\begin{itemize}[leftmargin=*]
  \item \code{ESPUI.number(\ldots)} --- поле ввода числа; \code{constrain()} не даёт ввести угол вне 1--360;
  \item \code{B\_UP} --- кнопка отпущена: одно действие на один щелчок;
  \item \code{requestDeg.exchange(0)} --- забрать накопленное и обнулить одной неделимой операцией.
\end{itemize}

\subsubsection{Шаг 2. Подключаем мотор}

Добавляем функции из основной программы урока: \code{applyPhase()}, \code{release()} и
\code{stepperRun()}. Теперь цель считается в шагах, а подпись показывает, где по мнению платы стоит вал.

\codecap{Шаг 2. Кнопка «Вправо» поворачивает вал}
\codeinput[\footnotesize]{l25_espui_stepper_s2}

\begin{itemize}[leftmargin=*]
  \item \code{target += deg * STEPS\_PER\_REV / 360} --- градусы из панели превращаются в шаги;
  \item \code{stepperRun()} вызывается в каждом проходе \code{loop()} и шагает не чаще, чем раз в \code{STEP\_US};
  \item подпись обновляется раз в 300 мс: слать в телефон каждый шаг незачем.
\end{itemize}

\subsubsection{Шаг 3. «Влево» и «В начало»}

Последний шаг добавляет две кнопки. «Влево» --- та же просьба с обратным знаком. «В начало» просит
повернуться на угол, равный текущему, но со знаком минус.

\codecap{Шаг 3. Три кнопки поворота}
\codeinput[\footnotesize]{l25_espui_stepper}

При включении плата считает положение нулевым, а настоящее положение вала она не знает: «Текущий угол»
всегда стартует с 0°.

\begin{trouble}
\begin{itemize}[leftmargin=*]
  \item \textbf{Двигатель жужжит и дрожит на месте} --- перепутан порядок проводов IN1--IN4 или шаги слишком частые.
  \item \textbf{Вал вращается не в ту сторону} --- поменяйте порядок выводов в массиве \code{PINS} на обратный.
  \item \textbf{Двигатель сильно греется} --- обмотки остаются под током на стоянке; вызывайте \code{release()}.
\end{itemize}
\end{trouble}

\begin{summary}
\begin{itemize}[leftmargin=*]
  \item Шаговый двигатель поворачивается одинаковыми шагами; положение --- это счётчик шагов.
  \item Порядок включения обмоток задаёт направление, пауза между шагами --- скорость.
  \item Двигатель подключают через драйвер и питают отдельно; на стоянке ток с обмоток снимают.
\end{itemize}
\end{summary}

\begin{selfcheck}
\begin{enumerate}
  \item Чем шаговый двигатель отличается от сервопривода и от обычного мотора?
  \item Сколько полушагов нужно сделать 28BYJ-48, чтобы повернуть вал на 45°?
  \item Что произойдёт, если подавать шаги слишком часто, и почему это опасно для программы?
\end{enumerate}
\end{selfcheck}

\begin{tasks}
\begin{enumerate}
\item \lv{1} Сделайте один оборот в одну сторону и один в обратную. Сколько шагов в обороте вашего мотора?
\item \lv{1} Запишите по датасету мотора ток обмотки и сверьте с током драйвера.
\item \lv{2} Добавьте команду \code{speed N}, которая меняет паузу между шагами.
\item \lv{2} Добавьте команду \code{angle N}: повернуться на N градусов.
\item \lv{3} Сделайте плавный разгон и торможение: первые и последние 200 шагов идут с увеличенной паузой.
\item \lv{4} Добавьте концевой выключатель: при включении мотор медленно вращается до него и принимает это положение за ноль.
\end{enumerate}
\end{tasks}
